%%%%%%%%%%%%%%%%%%%%%%%%%%%%%%%%%%%%%%%%%

%----------------------------------------------------------------------------------------
%	PACKAGES AND OTHER DOCUMENT CONFIGURATIONS
%----------------------------------------------------------------------------------------

\documentclass[11pt,fleqn]{book} % Default font size and left-justified equations

%----------------------------------------------------------------------------------------

%----------------------------------------------------------------------------------------
%	VARIOUS REQUIRED PACKAGES AND CONFIGURATIONS
%----------------------------------------------------------------------------------------

\NeedsTeXFormat{LaTeX2e}
\ProvidesClass{structure}[2019/02/04]
\LoadClassWithOptions{book}
% geometry
\RequirePackage{geometry}
\geometry{%
	left=1cm,%
	right=1cm,%
	top=1.5cm,%
	bottom=1.5cm,%
	bindingoffset=0.8cm}
% math font
\RequirePackage{amsmath}
\RequirePackage{amssymb}
\RequirePackage{amsthm}
\RequirePackage{mathtools}
\RequirePackage{newpxmath}

% Packages
\RequirePackage{lipsum} % Inserts dummy text
\RequirePackage{booktabs} % Required for nicer horizontal rules in tables
% language
\RequirePackage[no-math]{fontspec}% no math font
\setmainfont[Ligatures=TeX,AutoFakeBold=0.5,AutoFakeSlant=0.25]{Antioch Text}% roman font
\setsansfont[Ligatures=TeX,AutoFakeBold=0,AutoFakeSlant=0.25]{Khmer M1}% sans serif font
\setmonofont[Ligatures=TeX,AutoFakeBold=0.5,AutoFakeSlant=0.25]{Angkor}% typewriter font
\XeTeXlinebreaklocale "km"% line break rule
\XeTeXlinebreakskip = 0pt plus 0.50pt minus 0.25pt% line break skip
\newcommand{\kv}{\fontspec[Scale = 0.8333333334, Script=Khmer]{Khmer Vanara}\selectfont}
\newcommand{\km}{\fontspec[Scale = 0.8333333334, Script=Khmer]{Khmer M2}\selectfont}
\newcommand{\ak}{\fontspec[Scale = 0.8333333334, Script=Khmer]{Angkor}\selectfont}
\newcommand{\kml}{\fontspec[Scale = 0.8333333334, Script=Khmer]{Khmer M1}\selectfont}
\newcommand{\kos}{\fontspec[Scale = 0.8333333334, Script=Khmer]{Khmer OS System}\selectfont}
\newcommand{\kbk}{\fontspec[Scale = 0.8333333334, Script=Khmer]{Khmer OS Bokor}\selectfont}
\newcommand{\kb}{\fontspec[Scale = 0.8333333334, Script=Khmer]{Khmer OS Battambang}\selectfont}
\newcommand{\mc}{\fontspec{Monotype Corsiva}\selectfont}
%\newcommand{\en}{\fontspec{Times New Roman}\selectfont}
\newcommand{\en}{\fontspec{Abyssinica SIL}\selectfont}
\newcommand{\enb}{\fontspec[Script=Khmer]{Quicksand Medium}}
% khmer number
\def\khmer#1{\expandafter\@khmer\csname c@#1\endcsname}
\def\@khmer#1{\expandafter\@@khmer\number#1\@nil}
\def\@@khmer#1{%
	\ifx#1\@nil% terminate when encounter @nil
	\else%
	\ifcase#1 ០\or ១\or ២\or ៣\or ៤\or ៥\or ៦\or ៧\or ៨\or ៩\fi%
	\expandafter\@@khmer% recursively map the following characters
	\fi}
\def\khmernumeral#1{\@@khmer#1\@nil}
% khmer alphabet
\def\alpkh#1{\expandafter\@alpkh\csname c@#1\endcsname}
\def\@alpkh#1{%
	\ifcase#1% zero -> none
	\or ក\or ខ\or គ\or ឃ\or ង%
	\or ច\or ឆ\or ជ\or ឈ\or ញ%
	\or ដ\or ឋ\or ឌ\or ឍ\or ណ%
	\or ត\or ថ\or ទ\or ធ\or ន%
	\or ប\or ផ\or ព\or ភ\or ម%
	\or យ\or រ\or ល\or វ\or ស%
	\or ហ\or ឡ\or អ%
	\else%
	\@ctrerr% otherwise, counter error!
	\fi}
%
\def\khmershortdate{%
	\ifnum\day<10\@khmer{0\day}\else\@khmer\day\fi%
	/\@khmer\month%
	/\@khmer\year}
\def\khmerdate{%
	\ifnum\day<10\@khmer{0\day}\else\@khmer\day\fi~%
	\ifcase\month\or%
	មករា\or%
	កុម្ភៈ\or%
	មិនា\or%
	មេសា\or%
	ឧសភា\or%
	មិថុនា\or%
	កក្កដា\or%
	សីហា\or%
	កញ្ញា\or%
	តុលា\or%
	វិច្ឆិកា\or%
	ធ្នូ\fi~%
	\@khmer\year}
% color
\RequirePackage{xcolor} % Required for specifying colors by name
\definecolor{ocre}{RGB}{243,102,25} % Define the orange color used for highlighting throughout the book
% floating object
\RequirePackage{float}
% graphic
\RequirePackage{graphicx}
% table
\RequirePackage{colortbl}
\arrayrulecolor{header}
% picture
\RequirePackage{tikz}
\usetikzlibrary{shapes}
\usetikzlibrary{calc}
\usetikzlibrary{angles}
\usetikzlibrary{quotes}
\usetikzlibrary{arrows.meta}
\tikzset{%
	every picture/.style={%
		line width=1pt,%
		fill=header,%
		draw=header,%
		text=blue,%
	},%
	>={Stealth[length=6pt,inset=2pt]}%
}
%
\RequirePackage[many]{tcolorbox}

%%%%%%%%%%%%%%%%%%%%%%%%%%%%%%%%%%%%%%%%%%%%%%%%%%%%%%%%%%%%
% section spacing
\long\def\section{%
	\@startsection{section}{1}{\z@}{%
		-1.5ex\@plus-.2ex\@minus-.2ex}{%
		1.5ex\@plus.2ex\@minus.2ex}{\normalfont\bfseries}}
\long\def\subsection{%
	\@startsection{subsection}{2}{\z@}{%
		-1.5ex\@plus-.2ex\@minus-.2ex}{%
		1.5ex\@plus.2ex\@minus.2ex}{\normalfont\bfseries}}
\long\def\subsubsection{%
	\@startsection{subsubsection}{3}{\z@}{%
		-1.5ex\@plus-.2ex\@minus-.2ex}{%
		1.5ex\@plus.2ex\@minus.2ex}{\normalfont\bfseries}}
\long\def\paragraph{%
	\@startsection{paragraph}{4}{\z@}{%
		1.5ex\@plus.2ex\@minus.2ex}{-1em}{\normalfont\bfseries}}
\long\def\subparagraph{%
	\@startsection{subparagraph}{4}{\parindent}{%
		1.5ex\@plus.2ex\@minus.2ex}{-1em}{\normalfont\bfseries}}
%%%%%%%%%%%%%%%%%%%%%%%%%%%%%%%%%%%%%%%%%%%%%%%%%%%%%%%%%%%%

% enumerate
\RequirePackage[inline]{enumitem}
\RequirePackage{multicol}
\RequirePackage{pifont}
% set enumerate and itemize key for number of columns
% \RequirePackage{multicol}
\SetEnumitemKey{2}{before=\raggedcolumns\setlength{\multicolsep}{\topsep}\begin{multicols}{2},after=\end{multicols}}
\SetEnumitemKey{3}{before=\raggedcolumns\setlength{\multicolsep}{\topsep}\begin{multicols}{3},after=\end{multicols}}
\SetEnumitemKey{4}{before=\raggedcolumns\setlength{\multicolsep}{\topsep}\begin{multicols}{4},after=\end{multicols}}
\SetEnumitemKey{5}{before=\raggedcolumns\setlength{\multicolsep}{\topsep}\begin{multicols}{5},after=\end{multicols}}
\SetEnumitemKey{6}{before=\raggedcolumns\setlength{\multicolsep}{\topsep}\begin{multicols}{6},after=\end{multicols}}
%% set enumerate and itemize key for numerals
%\AddEnumerateCounter{\alpkh}{\@alpkh}{\char"1788}
%\AddEnumerateCounter{\khmer}{\@khmer}{\char"17e3}
\AddEnumerateCounter{\alpkh}{\@alpkh}{ឈ}
\AddEnumerateCounter{\khmer}{\@khmer}{៣}
% enumerate keys
\SetEnumitemKey{1}{label=\protect\color{ocre}\protect\sffamily\arabic*.}
\SetEnumitemKey{a}{label=\protect\color{ocre}\protect\sffamily\alph*.}
\SetEnumitemKey{A}{label=\protect\color{ocre}\protect\sffamily\Alph*.}
\SetEnumitemKey{i}{label=\protect\color{ocre}\protect\sffamily\roman*.}
\SetEnumitemKey{I}{label=\protect\color{ocre}\protect\sffamily\Roman*.}
\SetEnumitemKey{k}{label=\protect\color{ocre}\protect\sffamily\alpkh*.}
\SetEnumitemKey{m}{label=\protect\color{ocre}\protect\sffamily\khmer*.}
% itemize keys
\SetEnumitemKey{b}{label=\protect\color{ocre}\protect\ensuremath{\bullet}}
\SetEnumitemKey{c}{label=\protect\color{ocre}\protect\ensuremath{\circ}}
\SetEnumitemKey{d}{label=\protect\color{ocre}\protect\ensuremath{\diamond}}
\SetEnumitemKey{t}{label=\protect\color{ocre}\protect\ensuremath{\triangleright}}
\SetEnumitemKey{p}{label=\protect\color{ocre}\protect\ding{43}}
\SetEnumitemKey{w}{label=\protect\color{ocre}\protect\ding{45}}

% styling enumerate
\setlist[enumerate]{%
	parsep=0pt,%
	topsep=1ex,%
	partopsep=0pt,%
	leftmargin=*,%
	itemsep=.5ex,%
	labelsep=0.5ex}
\setlist[enumerate,1]{m}
\setlist[enumerate,2]{k}
\setlist[enumerate,3]{a}
\setlist[enumerate,4]{i}
% styling itemize
\setlist[itemize]{%
	parsep=0pt,%
	topsep=1ex,%
	partopsep=0pt,%
	leftmargin=*,%
	itemsep=.5ex,%
	labelsep=.5ex}
\setlist[itemize,1]{b}
\setlist[itemize,2]{c}
\setlist[itemize,3]{d}
\setlist[itemize,4]{t}
% sub enumerate
\RequirePackage{tasks}
\NewPatternFormat{k}{\@alpkh}% cntformats package
\NewPatternFormat{m}{\@khmer}% counter-format=tsk[l]
\settasks{%
	counter-format=tsk[k].,%
	label-format=\color{header}\sffamily,%
	label-offset=0.5ex,%
	label-align=right,%
	label-width=3ex,%
	item-indent=3.5ex,%
	column-sep=2ex,%
	before-skip=0ex,%
	after-skip=0ex,%
	after-item-skip=0.5ex%
}
\NewTasks[%
style=enumerate,%
counter-format=tsk[k].,%
label-format=\color{header}\sffamily,%
label-offset=0.5ex,%
label-align=right,%
label-width=3ex,%
item-indent=3.5ex,%
column-sep=2ex,%
before-skip=0ex,%
after-skip=0ex,%
after-item-skip=0.5ex]{Enumerate}[\item]
\NewTasks[%
style=itemize,
label=\ensuremath{\bullet},
label-format=\color{header},%
label-offset=0.5ex,%
label-align=right,%
label-width=1ex,%
item-indent=1.5ex,%
column-sep=2ex,%
after-item-skip=0.5ex]{Itemize}[\item]

%----------------------------------------------------------------------------------------
%	MAIN TABLE OF CONTENTS
%----------------------------------------------------------------------------------------

\RequirePackage{titletoc} % Required for manipulating the table of contents

% Part text styling
\titlecontents{part}[0cm]
{\addvspace{20pt}\centering\large\bfseries}
{}
{}
{}

% Chapter text styling
\titlecontents{chapter}[1.25cm] % Indentation
{\addvspace{12pt}\large\sffamily\bfseries} % Spacing and font options for chapters
{\color{ocre!60}\contentslabel[\Large\thecontentslabel]{1.25cm}\color{ocre}} % Chapter number
{\color{ocre}}  
{\color{ocre!60}\normalsize\;\titlerule*[.5pc]{.}\;\thecontentspage} % Page number

% Section text styling
\titlecontents{section}[1.25cm] % Indentation
{\addvspace{3pt}\sffamily\bfseries} % Spacing and font options for sections
{\contentslabel[\thecontentslabel]{1.25cm}} % Section number
{}
{\hfill\color{black}\thecontentspage} % Page number
[]

% Subsection text styling
\titlecontents{subsection}[1.25cm] % Indentation
{\addvspace{1pt}\sffamily\small} % Spacing and font options for subsections
{\contentslabel[\thecontentslabel]{1.25cm}} % Subsection number
{}
{\ \titlerule*[.5pc]{.}\;\thecontentspage} % Page number
[]

% List of figures
\titlecontents{figure}[0em]
{\addvspace{-5pt}\sffamily}
{\thecontentslabel\hspace*{1em}}
{}
{\ \titlerule*[.5pc]{.}\;\thecontentspage}
[]

% List of tables
\titlecontents{table}[0em]
{\addvspace{-5pt}\sffamily}
{\thecontentslabel\hspace*{1em}}
{}
{\ \titlerule*[.5pc]{.}\;\thecontentspage}
[]
 % Insert the commands.tex file which contains the majority of the structure behind the template

\begin{document}

%----------------------------------------------------------------------------------------
%	TITLE PAGE
%----------------------------------------------------------------------------------------

%\begingroup
%\thispagestyle{empty}
%\begin{tikzpicture}[remember picture,overlay]
%\coordinate [below=12cm] (midpoint) at (current page.north);
%\node at (current page.north west)
%{\begin{tikzpicture}[remember picture,overlay]
%\node[anchor=north west,inner sep=0pt] at (0,0) {\includegraphics[width=\paperwidth]{background}}; % Background image
%\draw[anchor=north] (midpoint) node [fill=ocre!30!white,fill opacity=0.6,text opacity=1,inner sep=1cm]{\Huge\centering\bfseries\sffamily\parbox[c][][t]{\paperwidth}{\centering The Search for a Title\\[15pt] % Book title
%{\Large A Profound Subtitle}\\[20pt] % Subtitle
%{\huge Dr. John Smith}}}; % Author name
%\end{tikzpicture}};
%\end{tikzpicture}
%\vfill
%\endgroup

%----------------------------------------------------------------------------------------
%	COPYRIGHT PAGE
%----------------------------------------------------------------------------------------

\newpage
~\vfill
\thispagestyle{empty}

\noindent Copyright \copyright\ 2013 John Smith\\ % Copyright notice

\noindent \textsc{Published by Publisher}\\ % Publisher

\noindent \textsc{book-website.com}\\ % URL

\noindent Licensed under the Creative Commons Attribution-NonCommercial 3.0 Unported License (the ``License''). You may not use this file except in compliance with the License. You may obtain a copy of the License at \url{http://creativecommons.org/licenses/by-nc/3.0}. Unless required by applicable law or agreed to in writing, software distributed under the License is distributed on an \textsc{``as is'' basis, without warranties or conditions of any kind}, either express or implied. See the License for the specific language governing permissions and limitations under the License.\\ % License information

\noindent \textit{First printing, March 2013} % Printing/edition date

%----------------------------------------------------------------------------------------
%	TABLE OF CONTENTS
%----------------------------------------------------------------------------------------

%\usechapterimagefalse % If you don't want to include a chapter image, use this to toggle images off - it can be enabled later with \usechapterimagetrue

\chapterimage{chapter_head_1.pdf} % Table of contents heading image

\pagestyle{empty} % No headers

\tableofcontents % Print the table of contents itself

\cleardoublepage % Forces the first chapter to start on an odd page so it's on the right

\pagestyle{fancy} % Print headers again

%----------------------------------------------------------------------------------------
%	PART
%----------------------------------------------------------------------------------------

\part{Part One}

%----------------------------------------------------------------------------------------
%	CHAPTER 1
%----------------------------------------------------------------------------------------

\chapterimage{chapter_head_2.pdf} % Chapter heading image

\chapter{Text Chapter}

\section{Paragraphs of Text}\index{Paragraphs of Text}

\lipsum[1-7] % Dummy text

%------------------------------------------------

\section{Citation}\index{Citation}

This statement requires citation \cite{book_key}; this one is more specific \cite[122]{article_key}.

%------------------------------------------------

\section{Lists}\index{Lists}

Lists are useful to present information in a concise and/or ordered way\footnote{Footnote example...}.

\subsection{Numbered List}\index{Lists!Numbered List}

\begin{enumerate}
\item The first item
\item The second item
\item The third item
\end{enumerate}

\subsection{Bullet Points}\index{Lists!Bullet Points}

\begin{itemize}
\item The first item
\item The second item
\item The third item
\end{itemize}

\subsection{Descriptions and Definitions}\index{Lists!Descriptions and Definitions}

\begin{description}
\item[Name] Description
\item[Word] Definition
\item[Comment] Elaboration
\end{description}

%----------------------------------------------------------------------------------------
%	CHAPTER 2
%----------------------------------------------------------------------------------------

\chapter{In-text Elements}

\section{Theorems}\index{Theorems}

This is an example of theorems.

\subsection{Several equations}\index{Theorems!Several Equations}
This is a theorem consisting of several equations.

\begin{theorem}[Name of the theorem]
In $E=\mathbb{R}^n$ all norms are equivalent. It has the properties:
\begin{align}
& \big| ||\mathbf{x}|| - ||\mathbf{y}|| \big|\leq || \mathbf{x}- \mathbf{y}||\\
&  ||\sum_{i=1}^n\mathbf{x}_i||\leq \sum_{i=1}^n||\mathbf{x}_i||\quad\text{where $n$ is a finite integer}
\end{align}
\end{theorem}

\subsection{Single Line}\index{Theorems!Single Line}
This is a theorem consisting of just one line.

\begin{theorem}
A set $\mathcal{D}(G)$ in dense in $L^2(G)$, $|\cdot|_0$. 
\end{theorem}

%------------------------------------------------

\section{Definitions}\index{Definitions}

This is an example of a definition. A definition could be mathematical or it could define a concept.

\begin{definition}[Definition name]
Given a vector space $E$, a norm on $E$ is an application, denoted $||\cdot||$, $E$ in $\mathbb{R}^+=[0,+\infty[$ such that:
\begin{align}
& ||\mathbf{x}||=0\ \Rightarrow\ \mathbf{x}=\mathbf{0}\\
& ||\lambda \mathbf{x}||=|\lambda|\cdot ||\mathbf{x}||\\
& ||\mathbf{x}+\mathbf{y}||\leq ||\mathbf{x}||+||\mathbf{y}||
\end{align}
\end{definition}

%------------------------------------------------

\section{Notations}\index{Notations}

\begin{notation}
Given an open subset $G$ of $\mathbb{R}^n$, the set of functions $\varphi$ are:
\begin{enumerate}
\item Bounded support $G$;
\item Infinitely differentiable;
\end{enumerate}
a vector space is denoted by $\mathcal{D}(G)$. 
\end{notation}

%------------------------------------------------

\section{Remarks}\index{Remarks}

This is an example of a remark.

\begin{remark}
The concepts presented here are now in conventional employment in mathematics. Vector spaces are taken over the field $\mathbb{K}=\mathbb{R}$, however, established properties are easily extended to $\mathbb{K}=\mathbb{C}$.
\end{remark}

%------------------------------------------------

\section{Corollaries}\index{Corollaries}

This is an example of a corollary.

\begin{corollary}[Corollary name]
The concepts presented here are now in conventional employment in mathematics. Vector spaces are taken over the field $\mathbb{K}=\mathbb{R}$, however, established properties are easily extended to $\mathbb{K}=\mathbb{C}$.
\end{corollary}

%------------------------------------------------

\section{Propositions}\index{Propositions}

This is an example of propositions.

\subsection{Several equations}\index{Propositions!Several Equations}

\begin{proposition}[Proposition name]
It has the properties:
\begin{align}
& \big| ||\mathbf{x}|| - ||\mathbf{y}|| \big|\leq || \mathbf{x}- \mathbf{y}||\\
&  ||\sum_{i=1}^n\mathbf{x}_i||\leq \sum_{i=1}^n||\mathbf{x}_i||\quad\text{where $n$ is a finite integer}
\end{align}
\end{proposition}

\subsection{Single Line}\index{Propositions!Single Line}

\begin{proposition} 
Let $f,g\in L^2(G)$; if $\forall \varphi\in\mathcal{D}(G)$, $(f,\varphi)_0=(g,\varphi)_0$ then $f = g$. 
\end{proposition}

%------------------------------------------------

\section{Examples}\index{Examples}

This is an example of examples.

\subsection{Equation and Text}\index{Examples!Equation and Text}

\begin{example}
Let $G=\{x\in\mathbb{R}^2:|x|<3\}$ and denoted by: $x^0=(1,1)$; consider the function:
\begin{equation}
f(x)=\left\{\begin{aligned} & \mathrm{e}^{|x|} & & \text{si $|x-x^0|\leq 1/2$}\\
& 0 & & \text{si $|x-x^0|> 1/2$}\end{aligned}\right.
\end{equation}
The function $f$ has bounded support, we can take $A=\{x\in\mathbb{R}^2:|x-x^0|\leq 1/2+\epsilon\}$ for all $\epsilon\in\intoo{0}{5/2-\sqrt{2}}$.
\end{example}

\subsection{Paragraph of Text}\index{Examples!Paragraph of Text}

\begin{example}[Example name]
\lipsum[2]
\end{example}

%------------------------------------------------

\section{Exercises}\index{Exercises}

This is an example of an exercise.

\begin{exercise}
This is a good place to ask a question to test learning progress or further cement ideas into students' minds.
\end{exercise}

%------------------------------------------------

\section{Problems}\index{Problems}

\begin{problem}
What is the average airspeed velocity of an unladen swallow?
\end{problem}

%------------------------------------------------

\section{Vocabulary}\index{Vocabulary}

Define a word to improve a students' vocabulary.

\begin{vocabulary}[Word]
Definition of word.
\end{vocabulary}

%----------------------------------------------------------------------------------------
%	PART
%----------------------------------------------------------------------------------------

\part{Part Two}

%----------------------------------------------------------------------------------------
%	CHAPTER 3
%----------------------------------------------------------------------------------------

\chapterimage{chapter_head_1.pdf} % Chapter heading image

\chapter{Presenting Information}

\section{Table}\index{Table}

\begin{table}[h]
\centering
\begin{tabular}{l l l}
\toprule
\textbf{Treatments} & \textbf{Response 1} & \textbf{Response 2}\\
\midrule
Treatment 1 & 0.0003262 & 0.562 \\
Treatment 2 & 0.0015681 & 0.910 \\
Treatment 3 & 0.0009271 & 0.296 \\
\bottomrule
\end{tabular}
\caption{Table caption}
\end{table}

%------------------------------------------------

\section{Figure}\index{Figure}

\begin{figure}[h]
\centering\includegraphics[scale=0.5]{placeholder}
\caption{Figure caption}
\end{figure}

%----------------------------------------------------------------------------------------
%	BIBLIOGRAPHY
%----------------------------------------------------------------------------------------

%\chapter*{Bibliography}
%\addcontentsline{toc}{chapter}{\textcolor{ocre}{Bibliography}}
%\section*{Books}
%\addcontentsline{toc}{section}{Books}
%\printbibliography[heading=bibempty,type=book]
%\section*{Articles}
%\addcontentsline{toc}{section}{Articles}
%\printbibliography[heading=bibempty,type=article]

%----------------------------------------------------------------------------------------
%	INDEX
%----------------------------------------------------------------------------------------

\cleardoublepage
\phantomsection
\setlength{\columnsep}{0.75cm}
\addcontentsline{toc}{chapter}{\textcolor{ocre}{Index}}
\printindex

%----------------------------------------------------------------------------------------

\end{document}
